\chapter{Аналитический раздел}

\section{Способы перехвата функций ядра}

\subsection{ftrace}

ftrace предоставляет возможности для трассировки функций. С его помощью можно отслеживать контекстные переключения, измерять время обработки прерываний, высчитывать время на активизацию заданий с высоким приоритетом и многое другое~\cite{stevens}.

Ftrace был разработан Стивеном Росгедтом и добавлен в ядро в 2008 году, начиная с версии 2.6.27. Ftrace — фреймворк, предоставляющий отладочный кольцевой буфер для записи данных. Собирают эти данные встроенные в ядро программы–трассировщики~\cite{stevens}.

Работает ftrace на базе файловой системы debugfs, которая в большинстве современных дистрибутивов Linux смонтирована по умолчанию.

Каждую перехватываемую функцию можно описать следующей структурой:

\begin{lstlisting}[language=C, frame=single, numbers=left, caption={Структура ftrace\_book}]
struct ftrace_book {
	const char *name;
	void *function;
	void *original;
	
	unsigned long address;
	struct ftrace_ops ops;
};
\end{lstlisting}


Поля структуры:
\begin{itemize}
	\item[---] name — имя перехватываемой функции;
	\item[---] function — адрес функции–сбергки, которая будет вызываться вместо перехваченной функции;
	\item[---] original — указатель на место, куда следует записать адрес перехватываемой функции, заполняется при установке;
	\item[---] address — адрес перехватываемой функции, заполняется при установке;
	\item[---] ops — служебная информация ftrace.
\end{itemize}

\begin{lstlisting}[language=C, frame=single, numbers=left, caption={Пример заполнения структуры ftrace\_book}]
#define HOOK(_name, _function, _original) \
{ \
	.name = (_name), \
	.function = (_function), \
	.original = (_original), \
}

static struct ftrace_book hooked_functions[] = {
	HOOK("sys_clone",  fh_sys_clone,  kreal_sys_clone),
	HOOK("sys_execve",  fh_sys_execve,  kreal_sys_execve),
};
\end{lstlisting}

\subsection{kprobes}

Kprobes — специализированное API, в первую очередь предназначенное для отладки и трассирования ядра. Этот интерфейс позволяет устанавливать пред- и постобработчики для любой инструкции в ядре, а также обработчики на вход и возврат из функции. Обработчики получают доступ к регистрам и могут их изменять.

K probes реализуются с помощью точек останова (инструкции int3), внедряемых в исполнимый код ядра, что позволяет устанавливать к probes в любом месте любой функции, если оно известно. Аналогично, kretprobes реализуются через подмену адреса возврата на стеке и позволяют перехватить возврат из любой функции.

При использовании k probes для получения аргументов функции или значений локальных переменных надо знать, в каких регистрах или где на стеке они лежат, и самостоятельно их оттуда извлекать. Для решения данной проблемы существует j probes — надстройка над k probes, самостоятельно извлекающая аргументы функции из регистров или стека и вызывающая обработчик, который должен иметь ту же сигнатуру, что и перехватываемые функция. Однако j probes объявлен устаревшим и удален из современных ядер.

\subsection{Linux Security API}

Linux Security API — интерфейс, созданный для перехвата функций ядра. В критических местах кода ядра расположены вызовы security–функций, которые в свою очередь вызывают коллёски, установленные security–модулем. Security–модуль может анализировать контекст операции и принимать решение о ее разрешении или запрете.

Для Linux Security API характерны следующие ограничения:

\begin{itemize}
	\item[---] security–модули не могут быть загружены динамически, они являются частью ядра и требуют его перекомпиляции;
	\item[---] в системе может быть только один security–модуль.
\end{itemize}

Таким образом, для использования Security API необходимо поставлять собственную сборку ядра, а также интегрировать дополнительный модуль с SELinux или AppArmor, которые используются популярными дистрибутивами.

\subsection{Модификация таблицы системных вызовов}

В ядре Linux все обработчики системных вызовов хранятся в таблице sys\_call\_table~\cite{linux-kernel}. Подмена значений в этой таблице приводит к смене поведения всей системы. Таким образом, сохранив старое значение обработчика и подставить в таблицу собственный обработчик, можно перехватить системный вызов.

Однако данный подход обладает следующими недостатками:

\begin{itemize}
	\item[---] Техническая сложность реализации, заключающаяся в необходимости обхода защиты от модификации таблицы, атомарное и безопасное выполнение замены.
	\item[---] Невозможность перехвата некоторых обработчиков. В ядрах до версии 4.16 обработка системных вызовов для архитектуры x86\_64 содержала целый ряд оптимизаций. Некоторые из них требовали того, что обработчик системного вызова являлся специальным переходником, реализованным на ассемблере. Соответственно, подобные обработчики порой сложно, а иногда и вовсе невозможно заменить на собственные, написанные на Си~\cite{habr-hooking}.
\end{itemize}

\subsection{Сплайсинг}

Сплайсинг заключается в замене инструкций в начале функции на безусловный переход, ведущий в обработчик. Оригинальные инструкции переносятся в другое место и исполняются перед переходом обратно в перехваченную функцию. Именно таким образом реализуется jump-оптимизация для kprobes. Используя сплайсинг, можно добиться тех же результатов, но без дополнительных расходов на kprobes и с полным контролем ситуации.

Сложность использования сплайсинга заключается в необходимости синхронизации установки и снятия перехвата, обхода защиты от модификации регионов памяти с кодом, инвалидации коней процессора после замены инструкций, дизассемблирования заменяемых инструкций и проверки на отсутствие переходов внутрь заменяемого кода.

\subsection{Сравнительный анализ способов перехвата}

Результаты сравнения различных способов перехвата функций ядра представлены в таблице 1.1

\begin{table}[h]
	\centering
	\caption{Результаты сравнения}
	\begin{tabular}{|p{3cm}|p{3cm}|p{3cm}|p{3cm}|p{3cm}|}
		\hline
		Способ перехвата & Необходимость перекомпиляции ядра & Возможность перехвата любых обработчиков & Доступ к аргументам функции через переменные & Необходимость обхода защиты от модификации регионов памяти \\
		\hline
		Linux Security API & + & --- & + & --- \\
		\hline
		Модификация таблиц системных вызовов & --- & --- & + & + \\
		\hline
		kprobes & --- & + & --- & --- \\
		\hline
		Сплайсинг & --- & + & + & + \\
		\hline
		ftrace & --- & + & + & --- \\
		\hline
	\end{tabular}
\end{table}

\section{Перехватываемые функции ядра}

\subsection{Функция getdents64}

Для сокрытия файла необходимо перехватить функцию ядра getdents64, так как она возвращает записи каталога.

\begin{lstlisting}[language=C, frame=single, numbers=left, caption={Функция getdents64}]
#include <fcntl.h>

int getdents64(unsigned int fd, struct linux_dirent64 *dirp, unsigned int count);
\end{lstlisting}

Системный вызов getdents64 читает несколько структур linux\_dirent64 из каталога, на который указывает открытый файловый дескриптор fd, в буфер, указанный в dirp. В аргументе count задается размер этого буфера.

Структура linux\_dirent64 определена следующим образом:

\begin{lstlisting}[language=C, frame=single, numbers=left, caption={Структура linux\_dirent64}]
struct linux_dirent64 {
	ino64_t    d_ino;
	off64_t    d_off;
	unsigned short  d_reclen;
	unsigned char   d_type;
	char    d_name[];
};
\end{lstlisting}

В d\_ino указан номер inode. В d\_off задается расстояние от начала каталога до начала следующей linux\_dirent64. В d\_reclen указывается размер данной linux\_dirent64. В d\_name задается имя файла, в d\_type — тип файла.

Таким образом, перехват системного вызова getdents64 позволяет удалить запись из списка записей каталога. Перехват вызова осуществляется при помощи создания указателя на системный вызов getdents64.

\begin{lstlisting}[language=C, frame=single, numbers=left, caption={Указатель на getdents64}]
static asmlinkage long (*real_sys_getdents64)(const struct pt_regs *);
\end{lstlisting}

\subsubsection{Функция unlink}

Удаление записей из каталога производится с помощью функции unlink.

\begin{lstlisting}[language=C, frame=single, numbers=left, caption={Функция unlink}]
#include <unistd.h>

int unlink(const char *pathname);
\end{lstlisting}

Эта функция удаляет запись из файла каталога и уменьшает значение счетчика ссылок на файл pathname. Если на файл указывает несколько ссылок, то его содержимое будет через них по-прежнему доступно.

Таким образом, для того, чтобы запретить удаление файла, необходимо перехватить системный вызов unlink. Перехват вызова осуществляется при помощи создания указателя на системный вызов.

\begin{lstlisting}[language=C, frame=single, numbers=left, caption={Указатель на unlink}]
static asmlinkage long (*real_sys_unlinkat)(struct pt_regs *regs);
\end{lstlisting}

\subsubsection{Функции open, write, read}

Запись в файл осуществляется при помощи системного вызова write, чтение из файла — read.

\begin{lstlisting}[language=C, frame=single, numbers=left, caption={Функция write и read}]
#include <fcntl.h>

ssize_t write(int fd, const void *buf, size_t size);
ssize_t read(int fd, void *buf, size_t count);
\end{lstlisting}

Функции write и read работают не с именем файла, а с файловым дескриптором. Однако для того, чтобы определить, можно ли выполнить операцию с файлом, необходимо знать его имя, так как файлы, для которых необходимо запретить те или иные операции, хранятся в виде списка имен этих файлов. Тогда следует перехватить функцию open, которая позволяет получить имя файла.

Системный вызов open осуществляет открытие файла.

\begin{lstlisting}[language=C, frame=single, numbers=left, caption={Функция open}]
#include <fcntl.h>

int open(int dirfd, const char *pathname, int flags);
int open(int dirfd, const char *pathname, int flags, mode_t mode);
\end{lstlisting}

При перехвате open, если имя файла есть в списке, то необходимо сохранить идентификатор процесса, открывающего файл. Затем при перехвате write и read ориентироваться не на имя файла, а на идентификатор процесса и файловый дескриптор. Так можно будет однозначно определить, необходимо ли запретить чтение из файла или запись в него.

Таким образом, для того, чтобы запретить чтение из файла и запись в него, необходимо перехватить системные вызовы open, write и read. Перехват осуществляется при помощи создания указателей на системные вызовы.

\begin{lstlisting}[language=C, frame=single, numbers=left, caption={Указатели на open, write и read}]
static asmlinkage long (*real_sys_open)(struct pt_regs *regs);
static asmlinkage long (*real_sys_write)(struct pt_regs *regs);
static asmlinkage long (*real_sys_read)(struct pt_regs *regs);
\end{lstlisting}

\section{Вывод}

В результате сравнительного анализа выбран способ перехвата функций ядра — ftrace, так как он позволяет перехватывать любые функции ядра и не требует его перекомпиляции. Для сокрытия файла необходимо перехватить функцию getdents64, для запрета чтения из файла и запись в файл — функции open, read и write, удаления — unlink.

